% !TEX program = xelatex
\documentclass[11pt]{beamer}

\definecolor{BNUBlue}{RGB}{0, 78, 161}
\usecolortheme[named=BNUBlue]{structure}
\setbeamertemplate{itemize items}[circle]
\setbeamercolor{itemize item}{fg=BNUBlue}
\setbeamertemplate{enumerate item}{\color{BNUBlue}\insertenumlabel.}
\setbeamercolor{enumerate item}{fg=BNUBlue}

\usepackage{ctex}
\usepackage{amsmath}
\usepackage{amssymb}
\usepackage{braket}
\usepackage{graphicx}
% \graphicspath{{pics/}}
\usepackage{subfigure}
\setbeamertemplate{caption}[numbered]
\usefonttheme{serif}
\usepackage{caption}
\usepackage[english]{babel}
\usepackage{calligra}
\setbeamertemplate{navigation symbols}{}

\usepackage{etoolbox}
\makeatletter
\defbeamertemplate*{footnote}{nonumber}{
    \parindent 1em\noindent
    \raggedright
    \insertfootnotetext\par
}
\setbeamertemplate{footnote}[nonumber]
\makeatother

\usepackage{times}
\usetheme{CambridgeUS}
\usecolortheme{dolphin}

\pgfdeclareimage[height=0.6cm]{BNU001}{BNU001}
\logo{\pgfuseimage{BNU001}}

\usepackage{fontspec}
\setmainfont{Times New Roman}
\usepackage{xeCJK}
\setCJKmainfont{SimSun}[BoldFont= SimHei, ItalicFont = KaiTi]

\DeclareMathOperator{\Var}{Var}
\DeclareMathOperator{\Cov}{Cov}
\DeclareMathOperator{\E}{E}
\newcommand{\dd}{\mathrm{d}}
\newcommand{\Ga}{\mathrm{Gamma}}
\newcommand{\Be}{\mathrm{Beta}}
\newcommand{\Exp}{\mathrm{Exp}}

\AtBeginSection[]{
    \begin{frame}[shrink]
        \frametitle{本节内容}
        \footnotesize
        \tableofcontents[currentsection]
    \end{frame}
}

\title[学术汇报]
{学术汇报}

\author[Xiang Mingxu]
{\fontsize{8}{10}\selectfont
    讲解人：Xiang Mingxu
    \\[10pt]
    \textnormal{\href{mailto:202411079392@mail.bnu.edu.cn}}
    {Email: \textit{202411079392@mail.bnu.edu.cn}}}

\institute[BNU]
{\fontsize{6}{10}\selectfont
    北京师范大学文理学院统计系\\
    Department of Statistics, Faculty of Arts and Sciences\\
    Beijing Normal University}

\date[August 2026]
{\fontsize{8}{10}
    \selectfont August 2026}

\begin{document}

\logo{}
\begin{frame}
    \titlepage
    \begin{center}
        \vspace{-15pt}
        \includegraphics[width=0.3\linewidth]{BNU001.png}
    \end{center}
\end{frame}

\logo{\pgfuseimage{BNU001}}
\begin{frame}[shrink]
    \frametitle{目录}
    \footnotesize
    \tableofcontents[hideallsubsections]
\end{frame}
\logo{}

\section{数理统计基本概念}

\subsection{统计量的定义与性质}

\begin{frame}
    \frametitle{统计量的定义}
    \small
    \textbf{定义：}由样本算出的量称为统计量（statistic），即统计量是样本的函数，不依赖于任何未知参数。

    \pause
    \vspace{0.5em}
    \textbf{说明：}
    \begin{enumerate}
        \item 统计量只与样本有关，不能与未知参数有关
        \item 样本具有两重性（具体数/随机变量），统计量也具有两重性；正因为统计量可视为随机变量，才有概率分布，这是统计推断的依据
        \item 选用什么统计量取决于问题性质，好的统计量应尽可能集中样本中与所讨论问题有关的信息
    \end{enumerate}

    \pause
    \vspace{0.3em}
    例：设 $X \sim N(\mu,\sigma^2)$，$X_1,\ldots,X_n$ 为 i.i.d.\ 样本，则
    $\bar{X}$, $\sum X_i^2$ 是统计量；\\
    $\sum (X_i-\mu)$, $\sum X_i/\sigma^2$
    当 $\mu,\sigma^2$ 未知时不是统计量.
\end{frame}

\subsection{常用统计量}

\begin{frame}
    \frametitle{样本均值和样本方差}
    \small
    \textbf{样本均值：}$\displaystyle \bar{X} = \frac{1}{n}\sum_{i=1}^n X_i$

    \textbf{样本方差：}$\displaystyle S^2 = \frac{1}{n-1}\sum_{i=1}^n (X_i - \bar{X})^2$

    \textbf{样本标准差：}$S = \sqrt{S^2}$

    \pause
    \vspace{0.3em}
    \textbf{基本性质：}
    \begin{enumerate}
        \item $\sum_{i=1}^n (X_i - \bar{X}) = 0$
        \item 设 $Y_i = aX_i + b$（$a \neq 0$），则
            $\bar{Y} = a\bar{X} + b$，$S_Y^2 = a^2 S_X^2$
        \item 对任意常数 $c$，有
            $\sum_{i=1}^n (X_i - c)^2 \ge \sum_{i=1}^n (X_i - \bar{X})^2$，
            等号仅当 $c = \bar{X}$ 时成立
    \end{enumerate}
\end{frame}

\begin{frame}
    \frametitle{性质 3 的证明：偏差平方和的最小性}
    \small
    思路：加减 $\bar{X}$，平方展开即可。

    \begin{align*}
        \sum_{i=1}^n (X_i - c)^2 &= \sum_{i=1}^n \bigl[(X_i - \bar{X}) + (\bar{X} - c)\bigr]^2 \\
        &= \sum_{i=1}^n (X_i - \bar{X})^2 + 2(\bar{X} - c)\underbrace{\sum_{i=1}^n (X_i - \bar{X})}_{=0} + \sum_{i=1}^n (\bar{X} - c)^2 \\
        &= \sum_{i=1}^n (X_i - \bar{X})^2 + n(\bar{X} - c)^2 \\
        &\ge \sum_{i=1}^n (X_i - \bar{X})^2.
    \end{align*}

    \pause
    交叉项为零（性质1），$n(\bar{X} - c)^2 \ge 0$ 恒成立，等号仅当 $c = \bar{X}$。

    结论：在偏差平方和最小的意义下，样本均值 $\bar{X}$ 是总体均值的最优估计。
\end{frame}

\begin{frame}
    \frametitle{样本矩}
    \small
    \textbf{样本 $k$ 阶原点矩：}
    \[
        a_{n,k} = \frac{1}{n}\sum_{i=1}^n X_i^k, \quad k = 1,2,\ldots
    \]
    当 $k = 1$ 时，$a_{n,1} = \bar{X}$，即样本均值。

    \pause
    \vspace{0.3em}
    \textbf{样本 $k$ 阶中心矩：}
    \[
        m_{n,k} = \frac{1}{n}\sum_{i=1}^n (X_i - \bar{X})^k, \quad k = 1,2,\ldots
    \]
    当 $k = 2$ 时，$m_{n,2} = \frac{n-1}{n}S^2$。
\end{frame}

\begin{frame}
    \frametitle{样本相关系数：从余弦相似度看线性关系}
    \small
    记中心化向量：
    \[
        \tilde{\mathbf{x}} = (X_1-\bar{X},\ldots,X_n-\bar{X})^\top,\quad
        \tilde{\mathbf{y}} = (Y_1-\bar{Y},\ldots,Y_n-\bar{Y})^\top.
    \]

    两个向量的余弦相似度定义为：
    \[
        \cos\theta = \frac{\tilde{\mathbf{x}}^\top\tilde{\mathbf{y}}}{\|\tilde{\mathbf{x}}\|\cdot\|\tilde{\mathbf{y}}\|}
        = \frac{\sum (X_i-\bar{X})(Y_i-\bar{Y})}{\sqrt{\sum (X_i-\bar{X})^2 \cdot \sum (Y_i-\bar{Y})^2}} = r.
    \]
    这正是样本相关系数 $r$。

    分子：两个中心化向量的内积；分母：两向量模长（长度）的乘积。
    本质：$\tilde{\mathbf{x}}$ 与 $\tilde{\mathbf{y}}$ 夹角的余弦值。
\end{frame}


\end{document}
